\section{Methodology}

We designed an experiment to test whether coefficient of variation (CV) of linear probe accuracy trajectories can distinguish spurious from core features.

\subsection{Feature Extraction}

We use CLIP ViT-B/16 \citep{radford2021learning} as the feature extractor. CLIP was pretrained on 400 million image-text pairs using contrastive learning. For each Waterbirds image, we extract the 512-dimensional embedding from the final layer.

\subsection{CV Measurement Protocol}

For each concept (background, bird\_type), we train logistic regression probes with $C \in \{0.001, 0.01, 0.1, 1, 10, 100\}$. The regularization strength $C$ serves as an ``epoch proxy.''

To measure emergence uniformity, we train probes on 5 random 20\% subsets of training data (seed=42). For each concept:
\begin{enumerate}
    \item Train probes on each subset across all $C$ values
    \item Record accuracy at each (subset, $C$) combination
    \item Compute improvement rate: $\text{acc}[C_i] - \text{acc}[C_{i-1}]$
    \item Calculate CV of improvement rates across subsets
\end{enumerate}

\textbf{Hypothesis}: Spurious features (background) should show low CV (uniform emergence); core features (bird\_type) should show high CV (differential emergence).

\subsection{Classification Evaluation}

Waterbirds provides ground truth: background is spurious (95\% correlated with label), bird\_type is core. We evaluate with AUC for binary classification, with gate threshold AUC $\geq$ 0.75.

\subsection{Experimental Configuration}

\begin{table}[h]
\centering
\begin{tabular}{@{}ll@{}}
\toprule
Parameter & Value \\
\midrule
Feature extractor & CLIP ViT-B/16 \\
Embedding dim & 512 \\
Probe & LogisticRegression \\
C sweep & [0.001, 0.01, 0.1, 1, 10, 100] \\
n\_subsets & 5 \\
Dataset & Waterbirds (4795 train samples) \\
\bottomrule
\end{tabular}
\end{table}
